\nextpage
\section*{Выводы и рекомендации}
На учебном корпусе для готового сканера re2c получено 2,478~нс/байт. В контролируемом сравнении двух представлений одного автомата на том же корпусе branch имеет меньшую медиану, а превосходство table на трёх размерах смешанной строки сохраняется в каждом раунде. Следовательно, ответ на вопрос о скорости зависит и от реализации, и от того, какими путями текст проводит автомат между выдачами лексем.

\begin{enumerate}
\item \textbf{Для обработки учебного псевдокода начать с re2c.}
На данном корпусе его межквартильный интервал отделён от интервалов остальных трёх сканеров. Если вход в основном состоит из коротких имён, стоит отдельно измерить табличный вариант: на двухбуквенных именах он вышел на первое место.
\item \textbf{При выборе представления собственного автомата учитывать длину слов и пропусков.}
На именах из 2 и 8 букв выигрывала таблица, на именах из 32 и 128 букв~--- ветвления. Аналогичная смена обнаружена при укрупнении комментариев с сохранением их общего объёма. Число состояний автомата во всех запусках одинаково.
\item \textbf{Близкие варианты сравнивать по общим раундам и указывать разброс.}
На \code{abcde2} и на длинных комментариях два первых результата опыта~А практически совпали. Для собственной пары на учебном корпусе в трёх раундах из 16 порядок также оказался обратным большинству. Для выбора на другой машине требуется повторение серии.
\end{enumerate}

\subhead{Ограничения}
Все выводы относятся к одной машине, одной оптимизированной сборке, одному учебному корпусу и четырём семействам искусственных текстов. Перед замером проводился прогрев. Привязку к ядру и частоту процессора не фиксировали. Учитывались выдача лексем и контрольная сумма. Подготовка буферов, чтение с диска и время построения сканеров в оценку не входят. Аппаратные счётчики обращений к памяти и ошибок предсказания переходов не снимались. Межквартильный размах характеризует эту серию из 16 раундов.

\begingroup\small
\renewcommand{\refname}{Источники}
\begin{thebibliography}{9}\setlength{\itemsep}{2pt}\setlength{\parskip}{0pt}
\bibitem{ga} Новиков Ф.~А. Грамматики и автоматы. Материалы курса. СПб., 2026. Пп.~0.2.4, 1.2.1, 2.2.1, 2.6.1, 2.6.2, 3.2.1. Для слайдов в тексте приведены их номера внутри пункта.
\bibitem{flex} The flex project. Lexical Analysis With Flex. Разделы \href{https://westes.github.io/flex/manual/Matching.html}{«How the Input Is Matched»} и \href{https://westes.github.io/flex/manual/Options-for-Scanner-Speed-and-Size.html}{«Options for Scanner Speed and Size»} (ключ \code{-Cf}). \url{https://westes.github.io/flex/manual/}. Дата обращения 28.11.2026.
\bibitem{re2c} The re2c project. User manual (C/C++). Раздел «Options», параметр \code{-{}-code-model}, модель \code{goto-label}. \url{https://re2c.org/manual/manual_c.html}. Дата обращения 28.11.2026.
\end{thebibliography}
\endgroup
